\documentclass[10pt,a4paper]{amsart}
\usepackage[margin=19mm]{geometry}
\usepackage{amsmath,amssymb,mathtools}
\usepackage[scr=rsfs,cal=euler]{mathalfa}
\usepackage{xfrac}
\usepackage[unicode,hidelinks,bookmarks=false]{hyperref}
\newcommand{\ie}{i.e.\ }
\newcommand{\cf}{cf.\ }
\newcommand{\resp}{respectively\ }
\newcommand{\Subsection}{\subsection}
\providecommand{\nobreakdash}{\nobreak\hbox{-}}
\setlength{\emergencystretch}{2em}
\multlinegap0pt
\usepackage{amssymb}
\usepackage{xcolor}
\newtheorem{theorem}{Theorem}
 \def\sol{\text{Sol}_3}
\title{Gauss curvature solitons on invariant surfaces in the homogeneous space Sol}
\author{Rafael Belli, Rafael L\'opez}
\date{Source: arXiv:2605.17097v1 (CC BY 4.0)}
\begin{document}
\maketitle

\noindent\emph{Source and licence.} Gauss curvature solitons on invariant surfaces in the homogeneous space Sol. Version arXiv:2605.17097v1 (CC BY 4.0), distributed by arXiv under the Creative Commons Attribution 4.0 International licence, http://creativecommons.org/licenses/by/4.0/, as recorded in the arXiv licence metadata for this exact version and shown on the abstract page https://arxiv.org/abs/2605.17097v1. Full licence notice: \texttt{licenses/arxiv-2605.17097-CC-BY-4.0.txt}.\par\smallskip

\noindent\textit{Editorial scope.} Counted unit: the article's theorem t55 (source0.tex lines 602--610). The article's complete original proof of it is delivered with it (source0.tex lines 612--654, one \texttt{proof} environment of 43 lines). No mathematics is written, repaired or re-derived: the statement and the proof are the article's own lines, in the article's own notation, and no retained line is substituted or re-flowed.  Retained uncounted as the source-local context the proof needs: lines 237--242; lines 489--491; lines 509--515; lines 571--573; lines 575--577.  Nothing was omitted from any retained span, so the package carries no omission record.  No retained line contains a citation, so the wrapper carries no bibliography and no bibliography entry.  No retained line contains a figure environment, an image-inclusion command or drawing code, so no image is delivered and the package declares no asset dependency.  Editorial formatting only, in addition: an AMS \texttt{amsart} preamble that restates the article's own declarations and macros used by the retained lines, namely the environments \texttt{theorem} on separate counters, together with the standard packages those lines need.  The retained environments print in this document as Theorem 1 in order of appearance, whereas the article prints the same items with the numbering of its own full document.  The complete native source of the article as distributed in the arXiv e-print for this exact version is bundled unchanged as \texttt{sources/200766/source0.tex}.
\begin{equation}\label{1pp}
\begin{split}
y'(s) e^{-z(s)} &= \cos\theta(s), \\
z'(s) &= \sin\theta(s). 
\end{split}
\end{equation}

 \begin{equation}\label{121}
 \theta' \cos\theta+ \cos^2\theta = \sin\theta\, e^{-z}.
 \end{equation}

\begin{theorem}\label{t52}
Let $\Sigma$ be an $F_1$-invariant surface in $\sol$ generated by a curve $\gamma$. Suppose that the angle function is a constant, $\theta(s)=c$.
\begin{enumerate}
 \item The surface $\Sigma$ cannot be an extrinsic $F_2$-soliton.
 \item If $\Sigma$ is an intrinsic $F_2$-soliton, then $\Sigma$ is a horizontal plane $z = z_0$.
\end{enumerate}
\end{theorem}

\begin{equation}\label{211}
z'' - z'^2 - z' e^{-z} + 1 = 0.
\end{equation}

 \begin{equation}\label{eq211}
v v' - v^2 - v e^{-z} + 1 = 0, \quad \text{with } v \neq 0.
\end{equation} 

 \begin{theorem} \label{t55}
Let $\Sigma$ be an $F_1$-invariant extrinsic $F_2$-soliton generated by a curve $\gamma$ which is a solution of \eqref{1pp}--\eqref{121} with initial conditions $y(0)=y_0$, $z(0)=z_0$ and $\theta(0)=0$. Then $\gamma$ has the following geometric properties: 
\begin{enumerate}
 \item The curve $\gamma$ cannot be a horizontal line.
 \item The point $s=0$ is a strict global maximum for the height function $z(s)$. 
 \item For $s>0$, $\gamma$ ends at a finite arc-length and a finite height with vertical velocity ($\theta \to -\pi/2$).
 \item For $s<0$, $\gamma$ descends infinitely ($z \to -\infty$) and in the $(y,z)$ plane, $\gamma$ is asymptotic to a straight line of slope $ 1$.
\end{enumerate}
\end{theorem}

\begin{proof} 
As established in \eqref{eq211}, if $v=z'$, then $v=v(s)$ satisfies the second-order equation $v' = v^2 + v e^{-z} - 1$. Let us introduce the function $w(s) = e^{-z(s)} > 0$ and define a planar autonomous dynamical system by 
\begin{equation}\label{fase1}
\begin{cases}
v' = v^2 + vw - 1, \\
w' = -vw.
\end{cases}
\end{equation}
 
 The asymptotic behavior as $z \to +\infty$ corresponds to the limit $w \to 0$. The equilibrium points of the system are $P_1 = (1, 0)$ and $P_2 = (-1, 0)$. Since $v^2=1$, then $\theta(s)=\pm\frac{\pi}{2}$ and these points represent generating curves that are vertical lines. By computing the Jacobian matrix associated with \eqref{fase1}, it is not difficult to see that $P_1$ and $P_2$ are saddle points. 
 


Denote by $\eta(s)=(v(s),w(s))$ the trajectory associated with a solution $(y(s),z(s),\theta(s))$ of \eqref{1pp}--\eqref{121}. In such a case, $(v,w)\in \Theta:=[-1,1]\times (0,\infty)$. Note that $P_1$ and $P_2$ belong to the boundary of $\Theta$.
\begin{enumerate}
\item If $\gamma$ is a horizontal line, then $v = 0$ and thus $v' = 0$. However, the first equation of \eqref{fase1} yields a contradiction. The result is also a consequence of Theorem \ref{t52}.
 \item At $s=0$, we have $z'(0)=0$ and by \eqref{211}, $z''(0)=-1<0$. Thus $s=0$ is a strict local maximum. Since $z''<0$ for any critical point, no more local maximum can exist by the Rolle's theorem, making $s=0$ the strict global maximum.
 
 
\item For $s > 0$, $v(s)$ initially becomes negative because $v'(0)=-1$. As long as $v \in (-1, 0)$ and $w > 0$, the first equation of \eqref{fase1} gives $v' = (v^2 - 1) + vw<0$. Furthermore, since $w' = -vw > 0$, $w$ is strictly increasing, making $v'$ even more negative. Consequently, $v(s)$ strictly decreases without asymptotic limits in $(-1, 0)$. Since $v(s) \ge -1$, the trajectory $\eta$ must reach the boundary $v = -1$ of $\Theta$. Evaluating the vector field at this boundary gives $v'_{|v=-1} = -w < 0$. Since $\eta$ is strictly transverse to $v=-1$, it reaches this state in finite time $s_1 > 0$ at some finite value $w_1 > w_0$. Geometrically, this implies $\gamma$ cannot be smoothly extended further attaining a height $z_1 = -\log(w_1)$. Moreover, since $v'\to -1$, then $\gamma$ becomes vertical.


\item For $s < 0$, since $v(0)=0$ and $v'(0)=-1$, the function $v(s)$ must be strictly positive. Since $0<v\leq 1$, then the second equation of \eqref{fase1} yields $\frac{w'}{w}\geq -1$. Thus $w(s)\geq w(0)e^{-s}$. Letting $s\to-\infty$, we derive $\lim_{s\to-\infty}w(s)=+\infty$, which implies $z \to -\infty$. 

{\it Claim:} $\lim_{s \to -\infty} vw = 1$. Let us define $f(s) = v(s)w(s)$. Differentiating $f$ with respect to $s$ and using the system \eqref{fase1}, we obtain
$$f'=(z''-z'^2)e^{-z}=e^{-z}(f-1).$$
 Integrating from $s$ to $0$, with $s < 0$, we get:
\begin{equation*}
\begin{split}
 \log|f(0) - 1| - \log|f(s) - 1| &= \int_{s}^{0} w(\tau) \, d\tau \geq w(0) \int_{s}^{0} e^{-\tau} \, d\tau\\
 &=w(0)(e^{-s}-1).
 \end{split}
 \end{equation*}
Since $f(0) = 0$, and $h'<0$ for $f<1$, it follows that $f(s) < 1$ for all $s < 0$. Thus
$$ -\log(1 - f(s)) \geq w(0)(e^{-s}-1).$$
Letting $s\to-\infty$, we conclude that $\lim_{s \to -\infty} f(s) = 1$. This proves the claim. 

Finally, from the claim, $\lim_{s\to -\infty}v(s) = 0$ because $\lim_{s\to-\infty}w(s)=+\infty$. The slope of $\gamma$, viewed as the graph $z=z(y)$, is
$$ \frac{dz}{dy} = \frac{z'(s)}{y'(s)} = \frac{v(s)}{ \frac{1}{w(s)} \sqrt{1-v(s)^2}} = \frac{v(s)w(s)}{\sqrt{1-v(s)^2}} \to 1,\quad\mbox{as $s\to-\infty$}.$$
Consequently, the branch of $\gamma$ for $s<0$ descends infinitely while becoming asymptotic to a straight line of slope $1$.
 
\end{enumerate}
\end{proof}

\end{document}
